\section{Introduction}
\label{sec:intro}

AI assistants now perform many of the tasks a research project needs. They propose experiments, draft code, summarize literature, and turn an open-ended question into a confident answer in minutes. The question for researchers is no longer whether a model can produce research-shaped artifacts. It is whether the loop that produced them can be trusted as a research argument. A useful protocol for AI-assisted research should help a researcher take a vague idea and end up with a principled, evidence-backed deliverable that they can stand behind and that a reviewer can inspect.

Two failure modes keep contemporary AI-assisted research from clearing that bar. The first is post-hoc optimization. A loop tries many ideas, keeps the one that scores well on a metric, and writes up the result. The research argument is constructed after a metric is hit, not carried through from assumptions to evidence. Autonomous-scientist platforms \citep{lu2024aiscientist, gottweis2025aicoscientist} exemplify the pattern. The second is brittle correctness. Hallucinated citations, drifted scope, and confidently wrong derivations slip through because nothing in the loop separates an agent's self-talk from external observation, or a report from accepted evidence. Benchmarks for ML and science agents \citep{chan2024mlebench, chen2024scienceagentbench} document how widely the brittleness appears. Both pitfalls share a root. Freeform agent loops let activity outpace the structure a research argument requires.

Existing systems address adjacent layers of the same problem but leave the protocol layer empty. Autonomous-scientist platforms aim to remove the human from the loop, which moves the problem rather than solving it for researchers who still own the project. Component-level tools, including literature and idea agents \citep{baek2024researchagent, skarlinski2024paperqa} and execution frameworks \citep{yao2023react, wu2023autogen, jiang2025aide}, produce useful artifacts and execution patterns, but none of them defines what counts as evidence or who holds authority over the project's claims. Agent-native research artifacts \citep{liu2026ara} address reproducibility of finished work, leaving the upstream process where claims are formed unstructured. None of those responses defines the protocol that turns intent into a reviewable argument.

Atelier supplies the protocol layer those responses leave open. The system is project-local and file-backed. Every consequential exchange between the human researcher, agents, and external observations is recorded in durable files rather than in chat history. The human researcher remains principal throughout, setting the research direction, accepting evidence, and approving the project's claims. \Cref{sec:method} develops the design principles, the assumptions about how agent-assisted research actually fails, and the looped research lifecycle that organize the protocol.

By routing work and review through files, Atelier improves not wall-clock speed but the structure, auditability, and recoverability of research workflows. A project run through the protocol produces a reconstructible chain from intent to claim, with errors surfaced where they happen rather than left silent. Section 3 reports a pilot study where two sweeps probe these claims on a simulated research scenario.
